% CLAUDE REPORT::EFCAL-0001::2026-11-10_1350_PST
\documentclass[11pt]{article}
\usepackage[margin=0.9in]{geometry}
\usepackage{amsmath,amssymb,amsthm,hyperref}
\newtheorem{theorem}{Theorem}
\newtheorem{proposition}[theorem]{Proposition}
\newtheorem{corollary}[theorem]{Corollary}
\theoremstyle{remark}
\newtheorem{remark}[theorem]{Remark}
\title{EFCAL-0001: Where the Content of the Worst-Case Derivation Programme Sits\\
\large A conservation identity for B1/B3 and an Extended Frege calibration of $\overline W_{R_{\mathrm{mix341}}}$}
\author{Claude routine (supporting analysis; not an LSPV role region)}
\date{2026-11-10, 13:50 PST}
\begin{document}
\maketitle

\noindent\textbf{Status summary.}
Theorem~\ref{thm:cons}: \textbf{PROVED/EXACT} (propositional logic, every fixed $R$).
Corollary~\ref{cor:known}: \textbf{PROVED/EXACT} given the cited \textbf{KNOWN THEOREMS}.
Theorem~\ref{thm:ef}: \textbf{CONDITIONAL} on any completion of the historical
$R_{\mathrm{mix341}}$ grammar [P2-D1] having only the rule types listed below,
each polynomial-time checkable and locally sound. Its two simulation directions are
\textbf{KNOWN THEOREMS} applied to that model.
[P2-B1], [P2-B2]/B3, [P2-D1], [P2-C3] and P versus NP remain \textbf{OPEN}.
Nothing here is a lower bound.

\section{Shared objective}
For a fixed sound, complete calculus $R$ with unit-paid distinct derived nodes,
\[
W_R(F)=\min_{\pi:F\vdash_R\bot}\operatorname{cost}_R(\pi),\qquad
\overline W_R(N)=\max_{\substack{F\in\mathrm{UNSAT}\\ |\langle F\rangle|\le N}} W_R(F).
\]
Write $S(R)$ for ``$\overline W_R$ is not bounded by any polynomial''. The project's two load-bearing arrows are
B1$(R)$: $S(R)$; and B3$(R)$, the boxed bridge: $\mathrm{SAT}\in\mathrm P\Rightarrow \neg S(R)$.
The directive's premise holds and is kept: a free optimal derivation is a legitimate
\emph{lower-bound} relaxation, because $W_R$ is a minimum over all legal derivations, so
strategy search is never charged. This report asks a different question: how much of P versus NP does each arrow carry?

\section{Conservation identity}
\begin{theorem}[PROVED/EXACT]\label{thm:cons}
For every fixed $R$:
\begin{enumerate}
\item B1$(R)\wedge$B3$(R)\iff S(R)\wedge \mathrm{SAT}\notin\mathrm P$.
\item $\mathrm{SAT}\notin\mathrm P\Rightarrow$ B3$(R)$. So B3 is vacuously true if P$\neq$NP, and refuting B3 for any $R$ requires proving P$=$NP.
\item If $S(R)$ is a theorem, then B3$(R)\iff\mathrm{SAT}\notin\mathrm P$.
\end{enumerate}
\end{theorem}
\begin{proof}
B3$(R)\equiv \mathrm{SAT}\notin\mathrm P\vee\neg S(R)$. Conjoining with $S(R)$ gives (1) and (3), and (2) is immediate.
\texttt{research/agents/claude/efcal/verify\_calibration.py} checks all valuations (28 checks, PASS) as a guard against transcription errors only.
\end{proof}

\begin{corollary}[PROVED/EXACT from KNOWN THEOREMS]\label{cor:known}
Several systems have unconditional superpolynomial worst-case refutation lower bounds on CNF, i.e.\ $S(R)$ holds:
ordinary DAG Resolution (Haken 1985; Ben-Sasson--Wigderson 2001 for expander Tseitin),
Cutting Planes (Pudl\'ak 1997),
bounded-depth Frege (Ajtai 1988; Beame--Impagliazzo--Kraj\'\i\v{c}ek--Pitassi--Pudl\'ak--Woods 1992),
and Polynomial Calculus (Razborov 1998).
For each of them, the boxed bridge B3 is \emph{logically equivalent to} P$\neq$NP.
\end{corollary}

\begin{remark}[Consequence for the directive]
The ``weakest sufficient bridge'' is not weaker than the target whenever the calculus is weak enough for a lower bound to be known. The content of P$\neq$NP cannot be split away by choosing a weaker $R$. It sits entirely in B1 when $R$ is strong and entirely in B3 when $R$ is weak. This is the abstract form of the ``hostile controls'' already in the corpus: Haken PHP versus Cook's ER proof [P2-B3 evidence], expander Tseitin versus GF(2) rows [P2-D1 evidence, LSPV-3/464], and the clock calculus [P2-B6]. A candidate $R$ is worth pursuing only if \emph{neither} arrow is already known to be equivalent to P$\neq$NP or its negation. Per the world-ledger rule in \texttt{AGENTS.md}, a reformulation equivalent to P$\neq$NP is not progress.
\end{remark}

\section{Extended Frege calibration of the fixed mixed calculus}
\paragraph{Model.} Let $R$ be any completion of $R_{\mathrm{mix341}}$ whose paid rules are: binary Resolution; AND3 and LOCAL-CONVERT of bounded support; binary GF(2) row XOR; and the conservative signed two-input VARDEF ($x\leftrightarrow \ell_1\wedge\ell_2$ by three clauses). Original source clauses are free. Each rule instance must be polynomial-time checkable and locally sound. This matches the rule list in LSPV-1/357 and LSPV-1/406. [P2-D1] (authenticating the historical byte grammar) is \textbf{not} assumed solved. Exotic undeclared rules would void the hypothesis.

\begin{theorem}[CONDITIONAL on the model]\label{thm:ef}
There is a polynomial $p$ such that for every UNSAT CNF $F$ with $|\langle F\rangle|=N$,
\[
W_R(F)\le O\bigl(W_{\mathrm{ER}}(F)\bigr),\qquad
\mathrm{size}_{\mathrm{EF}}(F)\le p\bigl(N+W_R(F)\bigr).
\]
Hence $\overline W_R$ is polynomially bounded iff Extended Frege (equivalently Extended Resolution) is polynomially bounded on CNF. In particular:
\begin{itemize}
\item B1$(R_{\mathrm{mix341}})\iff$ EF is not polynomially bounded, i.e.\ a superpolynomial EF lower bound for some CNF family. This is the central open problem of propositional proof complexity.
\item B3$(R_{\mathrm{mix341}})\iff$ (P$=$NP $\Rightarrow$ EF polynomially bounded). This question is open and studied as such in the literature.
\end{itemize}
\end{theorem}
\begin{proof}[Proof sketch]
\emph{Lower simulation.} ER with extensions restricted to $x\leftrightarrow\ell_1\wedge\ell_2$ is p-equivalent to full ER, since every Boolean connective expands into a constant number of such gates. An ER refutation with $L$ lines maps line by line to RES and VARDEF nodes, using $O(L)$ paid nodes.

\emph{Upper simulation.} Every derived node is a clause, a GF(2) row, or a definition over at most $N+W_R(F)$ variables, so the certificate has $\mathrm{poly}(N+W_R)$ bits. Each rule's local soundness is a polynomial-time-checkable statement provable in PV/$S^1_2$; for a XOR row this is parity of a sum. Cook (1975) and Kraj\'\i\v{c}ek--Pudl\'ak (1990) show that EF p-simulates every proof system whose soundness is PV-provable, which bounds the EF size. Finally, EF and ER are p-equivalent on CNF (Cook--Reckhow 1979; Kraj\'\i\v{c}ek 1995).
\end{proof}

\paragraph{What this changes.}
\begin{enumerate}
\item The shared directive's ``fixed, sufficiently powerful logical system'' $R_{\mathrm{mix341}}$ is, up to polynomials, Extended Frege. A proof of [P2-B1] would therefore be a superpolynomial EF lower bound. No such bound is known for any family. Every hard family with a proved lower bound in the corpus is EF-easy:
\begin{itemize}
\item PHP has polynomial-size ER proofs (Cook 1976).
\item Tseitin formulas have polynomial-size Frege proofs via Gaussian elimination.
\item Random 3-CNF with $m\ge n^{1.4}$ clauses has polynomial-size TC$^0$-Frege refutations (M\"uller--Tzameret 2014).
\end{itemize}
Constant-density random 3-CNF remains an \textbf{OPEN} candidate for EF hardness only.
\item B3 for this $R$ is exactly the question Pich--Santhanam (ECCC TR23-199) attack. They give assumptions under which EF lower bounds imply P$\neq$NP; for example, that $S^1_2$ proves that a polynomial-time function generates anticheckers for SAT. The corpus already records two sufficient conditions: PV-provable rejection soundness of a polynomial-time decider (LSPV-1/357) and Pich--Santhanam anticheckers (LSPV-1/406). This report adds that, \emph{given} B1, these conditions are not just sufficient but must be strong enough to imply P$\neq$NP (Theorem~\ref{thm:cons}(1)).
\item The identity-system equivalence of LSPV-5/409 (B3 $\iff Q_A\le_W R$) specializes here to ``EF polynomially simulates the trivial identity system of every polynomial-time SAT decider''.
\end{enumerate}

\section{Barrier-transfer and literature preflight}
Main obstacles: (i) no superpolynomial lower bound for any EF-strength system; (ii) no unconditional transfer from proof-length lower bounds to algorithm lower bounds.

\begin{itemize}
\item \emph{Direct transfer:} Cook / Kraj\'\i\v{c}ek--Pudl\'ak reflection simulation, applied to the stated model.
\item \emph{Barrier:} Arteche--Khaniki--Pich--Santhanam (arXiv 2405.02232) note that proving unconditionally that ``a concrete proof system is not polynomially bounded $\Rightarrow$ NP$\not\subseteq$P/poly'' would itself give new circuit lower bounds. So any noncircular B3 must carry circuit-lower-bound strength.
\item \emph{Failed transfer:} NP-hardness of finding or approximating shortest proofs (Atserias--M\"uller 2019 for Resolution; corpus LSPV-1/470) bounds proof \emph{search}, not \emph{minimum length}, and the directive deliberately does not charge search.
\item \emph{Post-quantum ledger:} PQ items 1--7 are FAILED TRANSFER here; the issue is purely logical and proof-complexity-theoretic.
\item \emph{Knot/link handoff (2026-11-09):} no material link to B1/B3 found this run.
\end{itemize}

\section*{CORPUS REVIEWED}
\texttt{AGENTS.md}; root \texttt{proof1.txt} and \texttt{proof2.txt} (current \texttt{origin/main}, entries [P2-00], [P2-B1]--[P2-B15], [P2-D1], [P2-D2], latest [P2-L3-552]); \texttt{main.tex} ownership-marker survey; LSPV-1/357 (G\"odel reflection B3 firewall); LSPV-1/406 (feasible antichecker B3); LSPV-1/470 and /471 (shared-DAG optimization and backward shortcut); LSPV-5/409 (identity-system B3 boundary); LSPV-3/552 handoff; the 2026-11-09 knot/link handoff list; and the Claude world ledger (rows L13, L34--L40).

\section*{DUPLICATION / CONTRADICTION AUDIT}
\begin{itemize}
\item Sufficient-condition directions for B3 duplicate LSPV-1/357 and /406, and are credited there.
\item The B3 identity-system equivalence is LSPV-5/409.
\item New in this report: the conservation identity (Theorem~\ref{thm:cons}), its consequence for systems with known lower bounds (Corollary~\ref{cor:known}), and the two-sided EF equivalence of B1 (Theorem~\ref{thm:ef}). A search of \texttt{proof2.txt} found no entry stating that B1$(R_{\mathrm{mix341}})$ is equivalent to an EF lower bound.
\item No contradiction found.
\end{itemize}

\section*{WHAT CHANGED RELATIVE TO ROOT MAIN}
\texttt{main.tex} is unchanged. This role has no ownership region, so LSPV-5 should decide whether to absorb this result. One evidence line is appended to \texttt{proof2.txt}.

\section*{ADOPTED / REJECTED / RETIRED PRIOR COMPONENTS}
\begin{itemize}
\item Adopted: free optimal derivation as the cost measure; the $R_{\mathrm{mix341}}$ rule list; and the LSPV-5/409 boundary.
\item Rejected: treating B3 as a ``weaker'' goal than P$\neq$NP for any $R$ with a proved lower bound.
\item None retired.
\end{itemize}

\section*{STALE OR OVERSTATED CLAIMS FOUND}
None outright. Wording that calls the boxed bridge the ``weakest'' route should state that, together with B1, it is equivalent to $S(R)\wedge$P$\neq$NP. For $R_{\mathrm{mix341}}$, B1 alone is the EF lower-bound problem.

\section*{Next decisive targets}
\begin{enumerate}
\item \textbf{Red team:} find a rule in the authenticated [P2-D1] grammar that falls outside the model, for example unbounded-support LOCAL-CONVERT. Such a rule would break the upper simulation.
\item \textbf{Hypergraph role:} test constant-density random 3-CNF and Tseitin-on-expander variants against EF/ER, not only against Resolution and Res(lin).
\item \textbf{Lead:} state B3 as an explicit EF question and assess whether any PV-provably-sound polynomial SAT decider could exist without collapsing known results.
\end{enumerate}

\section*{References}
S.~Cook, Feasibly constructive proofs and the propositional calculus, STOC 1975.
S.~Cook and R.~Reckhow, JSL 44 (1979).
S.~Cook, SIGACT News 8(4) (1976).
J.~Kraj\'\i\v{c}ek and P.~Pudl\'ak, Quantified propositional calculi and fragments of bounded arithmetic, ZML 36 (1990).
J.~Kraj\'\i\v{c}ek, \emph{Bounded Arithmetic, Propositional Logic and Complexity Theory}, CUP 1995.
A.~Haken, TCS 39 (1985).
P.~Pudl\'ak, JSL 62 (1997).
A.~Razborov, Comput.\ Complexity 7 (1998).
S.~M\"uller and I.~Tzameret, APAL 165 (2014).
A.~Atserias and M.~M\"uller, FOCS 2019.
J.~Pich and R.~Santhanam, \emph{Towards P$\neq$NP from Extended Frege lower bounds}, ECCC TR23-199, arXiv:2312.08163, \url{https://eccc.weizmann.ac.il/report/2023/199/}.
N.~Arteche, E.~Khaniki, J.~Pich, R.~Santhanam, arXiv:2405.02232.

\medskip\noindent P versus NP remains \textbf{OPEN}.
\end{document}
% END_OF_VERIFIED_REPORT::CLAUDE::EFCAL-0001
